% Methods

\section{Method}
\label{sec:methods}

\subsection{Problem setup and notation}
\label{sec:problem}

We consider the two-dimensional incompressible Navier--Stokes equations on the
periodic torus $\T^2 = [0,2\pi]^2$, the classical periodic setting of
Leray~\cite{leray1934}:
\begin{equation}
  \partial_t u + (u\cdot\grad)u = \nu\,\Delta u + f,
  \qquad \grad\cdot u = 0,
  \label{eq:ns}
\end{equation}
where $u = (u_1,u_2): \T^2 \to \R^2$ is the velocity, $\nu > 0$ the kinematic
viscosity, and $f$ an externally prescribed body force. Throughout we use
single-mode Kolmogorov forcing
\begin{equation}
  f = F \sin y \; e_x, \qquad F > 0,
  \label{eq:forcing}
\end{equation}
whose curl is $\grad \times f = -F\cos y$. The Reynolds number is defined as
$\mathrm{Re} = 2U/\nu$, with $U$ the Kolmogorov base speed and $L = 2\pi$ the
domain size; equivalently $\mathrm{Re} = U/(\nu k)$ at the fundamental
$k = \pi/L = 1/2$. The pairing of $F$ and $\nu$ for each Reynolds number is
given per run in Table~\ref{tab:forced-params} (Section~\ref{sec:forced}).

Since the domain is periodic, the force has zero mean, and $\grad \cdot u = 0$,
there exists a zero-mean stream function $\psi$ with
\begin{equation}
  u = (\partial_y \psi, -\partial_x \psi),
  \qquad \omega := \grad \times u = -\Delta \psi,
  \label{eq:stream}
\end{equation}
where $\omega$ is the vorticity. This sign convention is used throughout.
Substituting \eqref{eq:stream} into \eqref{eq:ns} and taking the curl yields
the stream-function formulation
\begin{equation}
  \partial_t \psi
  = \nu\,\Delta \psi + \Delta^{-1} J(\psi, \Delta \psi) - F \cos y,
  \qquad
  J(\psi, \varphi) := \partial_x \psi\, \partial_y \varphi
                      - \partial_y \psi\, \partial_x \varphi,
  \label{eq:stream-ns}
\end{equation}
where $\Delta^{-1}$ is the exact Fourier inverse on the zero-mean subspace.
The velocity is recovered from $\psi$ via \eqref{eq:stream} at every step, so
the discrete velocity is divergence-free by construction (Invariant I1,
Section~\ref{sec:invariants}).

All fields are represented on an $N \times N$ tensor-product Fourier
collocation grid ($n = N^2$ grid points); $\Delta$ and $\Delta^{-1}$ are the
spectral operators, and all $L^2$ norms below are the discrete
ell-two norms, which are spectrally equivalent on the collocation grid.
Identifying a field with its $N \times N$ matrix of values (rows indexed by
$x$, columns by $y$), the Laplacian separates as
\begin{equation}
  \Delta = D_x \otimes I + I \otimes D_y,
  \qquad D_x := -\partial_{xx}, \quad D_y := -\partial_{yy},
  \label{eq:lap-sep}
\end{equation}
with $[D_x \otimes I, \; I \otimes D_y] = 0$. This separation is the
structural fact exploited by the scheme.

\subsection{Energy identity}
\label{sec:energy}

Let $E(t) = \tfrac{1}{2}\norm{u}_2^2 = \tfrac{1}{2}\norm{\grad\psi}_2^2$ be the
kinetic energy. Taking the $L^2$ inner product of \eqref{eq:stream-ns} with
$\Delta \psi$ (equivalently, pairing the vorticity equation with
$\Delta^{-1}\omega$) gives
\begin{equation}
  \frac{\dd E}{\dd t}
  = \underbrace{\inner{u}{f}}_{P_{\mathrm{in}}(t)}
  \;-\;
  \underbrace{\nu\, \norm{\Delta \psi}_2^2}_{P_{\mathrm{diss}}(t)},
  \label{eq:energy}
\end{equation}
where the advective term is exactly energy-neutral (a standard identity for
periodic, divergence-free flows; see Temam~\cite{temam1977}) and the forcing
power is
\begin{equation}
  P_{\mathrm{in}}(t)
  = \inner{u}{f}
  = F\,\inner{\psi_y}{\sin y}
  = -F\,\inner{\psi}{\cos y},
  \label{eq:pin}
\end{equation}
with the last equality by integration by parts on the torus. The forced
system possesses the Kolmogorov equilibrium
\begin{equation}
  \psi_K = -\frac{F}{\nu}\cos y,
  \qquad
  u_K = \frac{F}{\nu}\sin y \; e_x,
  \label{eq:kolmogorov}
\end{equation}
at which \eqref{eq:energy} gives $P_{\mathrm{in}} = P_{\mathrm{diss}}
= 2\pi^2 F^2/\nu$. That equilibrium is an exact solution of the discretised
equations for any $(F,\nu)$: the diffusion and forcing terms cancel, and the
advection of $\omega_K$ vanishes identically. It is not, however, the state the
runs are started from --- Section~\ref{sec:setup} records how far below it they
begin.
% [RESOLVED from solvers/ns_psi.py: the 2/3 mask is applied to the field-level
% nonlinear and advection terms and not to the viscous semigroup;
% energy_terms() computes the advection term explicitly rather than assuming it
% vanishes, with zeta = -A cos y so that -<psi,zeta> = +A<psi,cos y> = P_in.]

Two consequences matter for the validation design. First, under forcing the
kinetic energy is \emph{not} monotone: $E(t)$ relaxes toward the forced
statistical regime with $P_{\mathrm{in}}(t) \approx P_{\mathrm{diss}}(t)$ on
average, and the relevant forcing-aware invariant is developed by the theory
work (Invariant I2, Section~\ref{sec:invariants};
% [PENDING-THEORETICAL-RESEARCH: forcing-aware invariant, docs/theory/
% stability-error.md]).
Second, the unforced limit $F = 0$ reduces \eqref{eq:energy} to
$\dd E/\dd t = -P_{\mathrm{diss}} \le 0$, so the laminar Taylor--Green decay
test must exhibit monotone non-increasing $E(t)$ (Section~\ref{sec:setup}).

\subsection{SP-DLRA ansatz and the exact viscous step}
\label{sec:spd}

We seek $\psi$ in the rank-$r$ ansatz
\begin{equation}
  \Psi := \psi \approx U S V^{\top},
  \qquad U, V \in \R^{N \times r}, \quad S \in \R^{r \times r},
  \label{eq:ansatz}
\end{equation}
with orthonormal columns $U^{\top}U = V^{\top}V = I_r$ and $S \succeq 0$
symmetric ($\Psi = USV^{\top}$ is its thin SVD; $r \ll N$). The
\emph{structure-preserving} property we exploit is that the viscous flow
preserves this ansatz \emph{exactly}.

\begin{proposition}[Exact, rank-preserving viscous step]
\label{prop:viscous}
Let $\Psi = USV^{\top}$ have rank $r$ with full-rank factors, and let
$\Delta = D_x \otimes I + I \otimes D_y$ as in \eqref{eq:lap-sep}. Then for
all $t \ge 0$,
\begin{equation}
  e^{\nu t \Delta}\, \Psi
  =
  \left(e^{\nu t D_x} U\right)\, S \,
  \left(e^{\nu t D_y} V\right)^{\top}.
  \label{eq:viscous-exact}
\end{equation}
In particular the viscous flow preserves the rank-$r$ ansatz exactly: the
evolved factors are $\hat U = e^{\nu t D_x} U$ and
$\hat V = e^{\nu t D_y} V$, and the singular values $S$ are unchanged.
\end{proposition}

\begin{proof}
$D_x \otimes I$ and $I \otimes D_y$ commute, so
$e^{\nu t \Delta} = e^{\nu t (D_x \otimes I)}\, e^{\nu t (I \otimes D_y)}$.
For any matrix $X$, $e^{\nu t (D_x \otimes I)} X = (e^{\nu t D_x} X)$, since
$D_x \otimes I$ applies $D_x$ row-wise; and, since $D_y$ is symmetric,
$X\, e^{\nu t (I \otimes D_y)} = X\, (e^{\nu t D_y})^{\top}$, since
$I \otimes D_y$ applies $D_y$ column-wise. Applying both to
$\Psi = USV^{\top}$ gives \eqref{eq:viscous-exact}.
\end{proof}

\begin{remark}[Cost of the viscous step]
$e^{\nu t D_x} U$ is computed column-wise by one-dimensional FFTs in the $x$
direction at cost $O(r N \log N) = O(r n \log n)$ (likewise for $V$). The full
viscous step is therefore $O(r n \log n)$, exact, and requires no SVD. This
is the cost of the viscous stage. The rank rule is separate: as implemented,
it factorizes the whole field at four stage boundaries per step, so the
per-step cost is $\Theta(N^3)$ and \textbf{independent of $r$} --- $r=2$ and
$r=64$ cost the same, because the truncated reconstruction changes only which
columns of an already-computed factorization are used --- measured, the ratio of
the $r=64$ to the $r=2$ full-step cost is $1.13$ at $N=64$ and within $5\%$ at
the two finer grids, against per-configuration timing spreads of order $1\%$.
Per-stage rank updates would invert this; we report the implemented behaviour
rather than the asymptotic one.
\end{remark}

\begin{remark}[What is preserved]
Equation~\eqref{eq:viscous-exact} says that the viscous semigroup is
invariant under the ansatz: the approximation error of the viscous step is
\emph{zero}, not merely small. This is the sense in which the scheme is
structure-preserving with respect to the viscous part, in the spirit of
separable/projector-splitting low-rank methods
\cite{koch2007, einkemmer2023, cui2026, goutaudier2026}. The nonlinear
part is \emph{not} invariant under the ansatz and is handled by a projected
step (Section~\ref{sec:proj-step}).
\end{remark}

\subsection{The projected nonlinear step}
\label{sec:proj-step}

Let $\Psi^n = USV^{\top}$ be the rank-$r$ approximation at time $t_n$ and let
$\Delta t$ be the time step. One SP-DLRA step is:

\begin{enumerate}
  \item \textbf{Exact viscous half-step (midpoint state).}
  \begin{equation}
    \Psi_m := e^{\nu \Delta t\, \Delta / 2}\, \Psi^n
    = \hat U S \hat V^{\top},
    \qquad
    \hat U := e^{\nu \Delta t D_x/2}\, U,
    \quad
    \hat V := e^{\nu \Delta t D_y/2}\, V,
    \label{eq:midpoint}
  \end{equation}
  by Proposition~\ref{prop:viscous} with $t = \Delta t/2$.
  \item \textbf{Nonlinear residual on the full grid.}
  \begin{equation}
    N(\Psi_m) := \Delta^{-1} J(\Psi_m, \Delta \Psi_m) - F \cos y.
    \label{eq:nonlin}
  \end{equation}
  \item \textbf{Projected midpoint update.} With the projection
  $\Pi_{\hat U, \hat V} G = \hat U (\hat U^{\top} G) \hat V^{\top}$,
  \begin{equation}
    \Psi^{n+1}
    =
    e^{\nu \Delta t\, \Delta}\, \Psi^n
    + \Delta t \; \Pi_{\hat U, \hat V}\, N(\Psi_m).
    \label{eq:step}
  \end{equation}
  Since $e^{\nu \Delta t \Delta} \Psi^n = \hat U_{\Delta t} S
  \hat V_{\Delta t}^{\top}$ (Proposition~\ref{prop:viscous}) with
  $\hat U_{\Delta t} = e^{\nu \Delta t D_x} U$,
  $\hat V_{\Delta t} = e^{\nu \Delta t D_y} V$, step \eqref{eq:step} is
  equivalently the coefficient update
  \begin{equation}
    S \leftarrow S + \Delta t\, \hat U_{\Delta t}^{\top}\,
    N(\Psi_m)\, \hat V_{\Delta t},
    \label{eq:s-update}
  \end{equation}
  on the viscously evolved factors (the Galerkin/projected integrator of
  Haasdonk and Ohlberger~\cite{haasdonk2012}).
  \item \textbf{Cleanup.} A thin SVD of $\Psi^{n+1}$ re-orthonormalizes the
  factors (rank at most $2r$ before truncation); singular values below
  tolerance are dropped (decay) and the residual indicator below may extend
  the rank (growth; Section~\ref{sec:rank-adaptation}). The SVD is of the
  \emph{whole} field, and is the dominant per-step cost of the scheme as
  implemented (Section~\ref{sec:cost}).
\end{enumerate}

The step is a \emph{projected exponential midpoint method}: the viscous part
is integrated exactly (Proposition~\ref{prop:viscous}) and the nonlinear part
is advanced by a midpoint rule on the projected dynamics, in the style of the
robust low-rank integrators of Einkemmer et al.~\cite{einkemmer2023} and
Ceruti et al.~\cite{ceruti2024}. For fixed factors the step is second order
in $\Delta t$; the projection error is the rank deficiency of the current
subspace, which the rank adaptation below is designed to control.
% [RESOLVED (D10): eq:step projects with the midpoint-evolved factors
% \hat U = e^{nu dt D_x/2} U, which the code evolves before the augmented step
% uses them, and the second-order property is asserted by
% test_reduced_path_is_second_order_in_dt and test_bug_is_second_order.]

\subsection{Adaptive rank}
\label{sec:rank-adaptation}

\paragraph{Growth (incremental SVD).}
Let $G := \Delta t\, N(\Psi_m)$ be the \emph{unprojected} nonlinear update and
$\Pi = \Pi_{\hat U, \hat V}$. We measure the component of $G$ outside the
current subspace by the residual indicator
\begin{equation}
  \eta := \frac{\norm{G - \Pi G}_F}{\norm{G}_F}.
  \label{eq:indicator}
\end{equation}
If $\eta > \eta_{\mathrm{tol}}$, the dominant direction of $G - \Pi G$ is
added to the subspace via an incremental singular value decomposition of the
augmented factors (the rank-adaptation technique of
Haasdonk and Ohlberger~\cite{haasdonk2012}), and the coefficients are
re-solved. Growth is the only \emph{active} rank change in the scheme. In the
forced runs the indicator is triggered within the first steps, the retained rank
reaches the budget, and it stays there for the rest of the run
(Section~\ref{sec:res-rank}).

\paragraph{Decay.}
At the cleanup SVD, singular values with $\sigma_i < \sigma_{\mathrm{tol}}$
(absolute) or $\sigma_i/\sigma_1 < \sigma_{\mathrm{tol}}$ (relative) are
dropped. This is a deterministic, cheap rule requiring no additional
decompositions. The implemented rule is a \emph{relative} amplitude cutoff of
$10^{-10}$ on the singular-value spectrum, tested every five steps, with the
rank requested between $2$ and the budget $2\lfloor N/3\rfloor+1$
(Section~\ref{sec:setup}); the state is carried at its own numerical rank, which
for the forced initial condition is seventeen. In the unforced Taylor--Green test
the state is a single Fourier mode, so the rule holds the rank at one from the
outset rather than spend rank on a spectrum that is not present
(Section~\ref{sec:results}).
% [RESOLVED from the run records: parameters.dlra_relative_amplitude_cutoff,
% dlra_check_every, dlra_min_rank, dlra_max_rank; the L1 run holds a fixed rank
% of one, so no tolerance is exercised there.]

\subsection{Invariants}
\label{sec:invariants}

We validate the scheme against four invariants, per reviewer decision D3:

\begin{itemize}
  \item \textbf{I1 (exact divergence-freeness).} The velocity is recovered as
  $u = (\psi_y, -\psi_x)$ from the (low-rank) stream function, so
  $\grad \cdot u = 0$ \emph{identically in exact arithmetic}. The measured
  residual never exceeds $1.1\times 10^{-11}$ over every method, rank and
  Reynolds number we ran, which is seven orders of magnitude below the reduced
  solver's own error ($\sim 10^{-4}$) and therefore cannot account for it. The
  residual is the roundoff of the discrete spectral derivative pair, and grows as
  the operators' conditioning does: $2.6\times$ from $N=64$ to $N=128$,
  consistently across methods. Section~\ref{sec:res-div} reports the measured
  band and the population it is taken over.
  \item \textbf{I2 (energy).} Unforced: $E(t)$ is monotone non-increasing
  (verified numerically in the Taylor--Green test). Forced: $E$ obeys
  \eqref{eq:energy} with $P_{\mathrm{in}} = -F\inner{\psi}{\cos y}$ and
  $P_{\mathrm{diss}} = \nu \norm{\Delta \psi}_2^2$; the forcing-aware
  invariant characterizing the statistical regime is given by the theory
  work.
  % [PENDING-THEORETICAL-RESEARCH: forcing-aware invariant and any
  % monotone/conserved quantity for the forced dynamics,
  % docs/theory/stability-error.md]
  \item \textbf{I3 (rank economy).} The adaptive rank $r(t)$ against the static
  POD count $r_{\mathrm{POD}} = 16$ --- the rank the static baselines were
  \emph{given}, not a count we measure. Their bases are fitted on the
  fluctuations of the reference run, including the initial condition, at a
  $99\%$-energy threshold. Because that rank is the requested cap and is the same
  in every run, it does not vary with the dynamics, and we therefore report the
  rank economy as a comparison against a fixed rank rather than as a measured
  saving.
  \item \textbf{I4 (turbulent fidelity).} Kinetic-energy time series,
  kinetic-energy spectrum, and dissipation statistics against the full-grid
  spectral reference.
\end{itemize}

\subsection{Cost model}
\label{sec:cost}

Per step, with $n = N^2$ grid points and rank $r$:
\begin{center}
\begin{tabular}{lll}
  \toprule
  Component & Cost & Remark \\
  \midrule
  Viscous half- and full-step (Prop.~\ref{prop:viscous}) & $O(r\, n \log n)$
    & factor-wise 1D FFT exponentials \\
  Nonlinear residual $N(\Psi_m)$ & $O(n \log n)$
    & $\Delta$, $\Delta^{-1}$ by FFT; $J$ pointwise;
    \emph{rank-independent} \\
  Projection $\Pi_{\hat U, \hat V} N(\Psi_m)$ & $O(r n)$ & \\
  Cleanup: thin SVD of the whole field, four per step & $\Theta(N^3)$ &
    \emph{rank-independent} \\
  \bottomrule
\end{tabular}
\end{center}

\paragraph{Timing protocol.} Per-step cost is the median of seven or more
repeats after a discarded warm-up, taken under two accountings --- the complete
reduced step, including the full-grid nonlinear evaluation and all projections,
and the time inside the whole-field factorizations alone --- each expressed as a
ratio against the pinned-thread full-grid reference. Configurations are
interleaved rather than run to completion one after another: on a shared node the
per-configuration median is unreliable, while the ratio between two interleaved
configurations is not. That is why we report the spreads and quote the band to one
significant figure. Two further facts belong here rather than in a footnote. The
rank is \emph{held fixed} for the timing, so the cost is measured at a fixed rank
and not under adaptation. And the benchmark state is an unforced multi-mode decay
rather than the forced initial condition, because a fixed low rank on a forced
field grows without bound and makes the factorization fail --- a property of that
configuration rather than a cost result.

The full-grid spectral reference has per-step cost $O(n \log n)$ with the same
asymptotic form. We emphasize the honest consequence: because the nonlinear
residual is evaluated on the full grid, its $O(n \log n)$ cost is
\emph{independent of the rank}, and we make no claim of per-step speedup:
SP-DLRA is measured at $2.2$--$3.5\times$ the full-grid step
(Section~\ref{sec:res-cost}, Figure~\ref{fig:cost}). Nor is there a compensating
memory benefit: peak memory is $2.3$ MiB ($N=64$) to $4.0$ MiB ($N=128$)
\emph{above} the full-grid step, because the state is a full field plus its
factors plus the factorisation workspace. We therefore identify no end-to-end
benefit in the regimes we have measured, and the case for the method rests on its
structural guarantees and its accuracy, not on efficiency. Establishing a regime
where the rank stays small over a long span would require evidence beyond the
$200$-step horizons used here, and we do not have it.
% [PENDING-CODER: the per-step cost and peak memory of SP-DLRA against the
% pinned-thread full-grid reference are measured, per grid, in
% cost_retiming.json and peak_memory.json. The static-baseline family is NOT
% timed, so there is no measured baseline cost ratio to report; this stays open
% until the coder times it.]
